\documentclass[10pt,a4paper]{amsart}
\usepackage[margin=19mm]{geometry}
\usepackage{amsmath,amssymb,mathtools}
\usepackage[scr=rsfs,cal=euler]{mathalfa}
\usepackage{xfrac}
\usepackage[unicode,hidelinks,bookmarks=false]{hyperref}
\newcommand{\ie}{i.e.\ }
\newcommand{\cf}{cf.\ }
\newcommand{\resp}{respectively\ }
\newcommand{\Subsection}{\subsection}
\providecommand{\nobreakdash}{\nobreak\hbox{-}}
\setlength{\emergencystretch}{2em}
\multlinegap0pt
\usepackage{bm}
\usepackage{amssymb}
\usepackage{mathtools}
\usepackage[shortlabels]{enumitem}
\usepackage{xcolor}
\usepackage{tcolorbox}
\newtheorem{lemma}{Lemma}
\newtheorem{theorem}{Theorem}
\title{Critical Norm Profiles for Finite-Prime Composition Operators on the Hardy Space of Dirichlet Series}
\author{Xiang Fang, Feng Guo, Aman Mishra, P. Muthukumar}
\date{Source: arXiv:2608.26041v1 (CC BY 4.0)}
\begin{document}
\maketitle

\noindent\emph{Source and licence.} Critical Norm Profiles for Finite-Prime Composition Operators on the Hardy Space of Dirichlet Series. Version arXiv:2608.26041v1 (CC BY 4.0), distributed by arXiv under the Creative Commons Attribution 4.0 International licence, http://creativecommons.org/licenses/by/4.0/, as recorded in the arXiv licence metadata for this exact version and shown on the abstract page https://arxiv.org/abs/2608.26041v1. Full licence notice: \texttt{licenses/arxiv-2608.26041-CC-BY-4.0.txt}.\par\smallskip

\noindent\textit{Editorial scope.} Counted unit: the article's theorem thm:boundary-Hankel-reduction (source0.tex lines 1900--1915). The article's complete original proof of it is delivered with it (source0.tex lines 1917--2083, one \texttt{proof} environment of 167 lines). No mathematics is written, repaired or re-derived: the statement and the proof are the article's own lines, in the article's own notation, and no retained line is substituted or re-flowed.  Retained uncounted as the source-local context the proof needs: lines 1255--1264.  Nothing was omitted from any retained span, so the package carries no omission record.  No retained line contains a citation, so the wrapper carries no bibliography and no bibliography entry.  No retained line contains a figure environment, an image-inclusion command or drawing code, so no image is delivered and the package declares no asset dependency.  Editorial formatting only, in addition: an AMS \texttt{amsart} preamble that restates the article's own declarations and macros used by the retained lines, namely the environments \texttt{lemma}, \texttt{theorem} on separate counters, together with the standard packages those lines need.  The retained environments print in this document as Lemma 1, Theorem 1 in order of appearance, whereas the article prints the same items with the numbering of its own full document.  The complete native source of the article as distributed in the arXiv e-print for this exact version is bundled unchanged as \texttt{sources/208615/source0.tex}.
\begin{lemma}
\label{lem:fp-boundedness-H-rho}
For a given $\boldsymbol{\rho}\in B_d,$  the matrix $\mathcal{H}_{\boldsymbol{\rho}}$ defines a bounded  operator on $\ell^2(\mathbb{N}_{0}^{d}).$
Moreover,
\[
\left\|\mathcal{H}_{\boldsymbol{\rho}}\right\|
\leq
\frac{2}{1+\sqrt{1-R_{\boldsymbol{\rho}}^2}}.
\]
\end{lemma}

\begin{theorem}
\label{thm:boundary-Hankel-reduction}
For $\boldsymbol{\rho}\in B_d$, the operator  \(\mathcal{H}_{\boldsymbol{\rho}}\) is unitarily equivalent to $\widetilde{B}_{\boldsymbol{\rho}}\oplus \mathbf{0},$ where
\[
\widetilde{B}_{\boldsymbol{\rho}}
=
\left(
(-1)^{N+M}(N+M)!h_N(\boldsymbol{\rho})h_M(\boldsymbol{\rho})
\right)_{N,M\geq 0}.
\]
Moreover, if $\boldsymbol{\rho}\neq \mathbf{0}$ then
$B_{\boldsymbol{\rho}} =
D_{\boldsymbol{\rho}}H_{\frac{R_{\boldsymbol{\rho}}}{2}}D_{\boldsymbol{\rho}},$
where $B_{\boldsymbol\rho}$ is unitarily equivalent to  $\widetilde B_{\boldsymbol\rho}$ by a diagonal unitary operator.

\end{theorem}

\begin{proof}
For $\alpha\in\mathbb{N}_{0}^{d},$ set
\[
w_{\alpha}
=
\frac{1}{\alpha!}
\prod_{j=1}^{d}
\left(\frac{\rho_j}{2}\right)^{\alpha_j}.
\]
Then
\[
h_N(\boldsymbol{\rho})^{2}
=
\sum_{|\alpha|=N}w_{\alpha}^{2}.
\]
The case  \(R_{\boldsymbol{\rho}}=0\)  is immediate, so assume henceforth that $R_{\boldsymbol{\rho}}>0$. Then,  for every \(N\) there is at least one multi-index \(\alpha\) with $|\alpha|=N$ and $w_{\alpha}>0.$
Thus $h_N(\boldsymbol{\rho})>0.$ Define
$\omega_N\in\ell^{2}(\mathbb{N}_{0}^{d})$ by
\[
(\omega_N)_{\alpha}
=
\begin{cases}
\displaystyle \frac{w_{\alpha}}{h_N(\boldsymbol{\rho})}, & |\alpha|=N,\\[6pt]
0, & |\alpha|\neq N.
\end{cases}
\]
The vectors $\omega_0,\omega_1,\omega_2,\ldots$ are orthonormal because they have disjoint total-degree supports and each has norm \(1\). Let $\mathcal{M}_{\boldsymbol{\rho}}=\overline{\operatorname{span}}\{\omega_N:N\geq 0\}$.

For $|\alpha|=N$ and $|\beta|=M,$ the matrix entry of \(\mathcal{H}_{\boldsymbol{\rho}}\) can be written as
\[
(\mathcal{H}_{\boldsymbol{\rho}})_{\alpha,\beta}
=
(-1)^{N+M}(N+M)!w_{\alpha}w_{\beta}.
\]

Let $x=(x_{\beta})_{\beta\in\mathbb{N}_{0}^{d}}$ be finitely supported. For $|\alpha|=N,$ we have
\[
\begin{aligned}
(\mathcal{H}_{\boldsymbol{\rho}}x)_{\alpha}
&=
\sum_{\beta\in\mathbb{N}_{0}^{d}}
(\mathcal{H}_{\boldsymbol{\rho}})_{\alpha,\beta}x_{\beta}  \\
&=
(-1)^Nw_{\alpha}
\sum_{M=0}^{\infty}
(-1)^M(N+M)!
\sum_{|\beta|=M}w_{\beta}x_{\beta}.
\end{aligned}
\]
Consequently, for each fixed \(N\),
\[
\sum_{|\alpha|=N}(\mathcal{H}_{\boldsymbol{\rho}}x)_{\alpha}e_{\alpha}= (-1)^N
\sum_{M=0}^{\infty}
(-1)^M(N+M)!
\sum_{|\beta|=M}w_{\beta}x_{\beta} \sum_{|\alpha|=N}w_{\alpha}e_{\alpha}
\]
Since $\sum_{|\alpha|=N}w_{\alpha}e_{\alpha}
= h_N(\boldsymbol{\rho})\omega_N,$
\[
\mathcal{H}_{\boldsymbol{\rho}}x =  \sum_{N=0}^{\infty}\left (
\sum_{M=0}^{\infty}
(-1)^{(N+M)}(N+M)!
\sum_{|\beta|=M}w_{\beta}x_{\beta} \right  )h_N(\boldsymbol{\rho})\omega_N
\]
Therefore, $\mathcal{H}_{\boldsymbol{\rho}}x\in \mathcal{M}_{\boldsymbol{\rho}}$  for every finitely supported \(x\). Since \(\mathcal{H}_{\boldsymbol{\rho}}\) is bounded (from Lemma \ref{lem:fp-boundedness-H-rho}), it follows that
$\operatorname{Ran}\mathcal{H}_{\boldsymbol{\rho}}
\subset \mathcal{M}_{\boldsymbol{\rho}}.$


Now let $x\in\mathcal{M}_{\boldsymbol{\rho}}^{\perp}.$
Then, for every $M\in\mathbb{N}_{0},$ we have
\[
0
=
\left\langle x,\omega_M\right\rangle
=
\frac{1}{h_M(\boldsymbol{\rho})}
\sum_{|\beta|=M}x_{\beta}w_{\beta}.
\]
Hence
\[
\sum_{|\beta|=M}w_{\beta}x_{\beta}=0
\qquad
\text{for every}
\qquad
M\in\mathbb{N}_{0}.
\]
Applying the preceding formula  for $\mathcal{H}_{\boldsymbol{\rho}}$ first to finitely supported vectors and then using density, we obtain

\[
\mathcal{H}_{\boldsymbol{\rho}}x=0.
\]
Thus $\mathcal{H}_{\boldsymbol{\rho}}$ annihilates
$\mathcal{M}_{\boldsymbol{\rho}}^{\perp}.$

Define $U:\ell^2(\mathbb{N}_{0})\to \mathcal{M}_{\boldsymbol{\rho}}$  by
\[
Ue_N=\omega_N,
\qquad
N\in\mathbb{N}_{0}.
\]
Since the vectors $\omega_N$ form an orthonormal basis for
$\mathcal{M}_{\boldsymbol{\rho}},$ the map \(U\) is unitary.

We compute the matrix of $U^{*}\mathcal{H}_{\boldsymbol{\rho}}U.$ For $|\alpha|=N$ and
$M\in\mathbb{N}_{0},$ we have
\[
\begin{aligned}
(\mathcal{H}_{\boldsymbol{\rho}}\omega_M)_{\alpha}=\sum_{\beta\in\mathbb{N}_0^d}(\mathcal{H}_{\boldsymbol{\rho}})_{\alpha,\beta}(\omega_M)_{\beta}
&=
\sum_{|\beta|=M}
(-1)^{N+M}(N+M)!w_{\alpha}w_{\beta}
\frac{w_{\beta}}{h_M(\boldsymbol{\rho})}                                      \\
&=
(-1)^{N+M}(N+M)!w_{\alpha}
\frac{1}{h_M(\boldsymbol{\rho})}
\sum_{|\beta|=M}w_{\beta}^{2}                              \\
&=
(-1)^{N+M}(N+M)!w_{\alpha}h_M(\boldsymbol{\rho}).
\end{aligned}
\]
Since $w_{\alpha}=h_N(\boldsymbol{\rho})(\omega_N)_{\alpha}$ for every coordinate $\alpha$ with $|\alpha|=N$,
\[
(\mathcal{H}_{\boldsymbol{\rho}}\omega_M)_{\alpha}= (-1)^{N+M}(N+M)!h_N(\boldsymbol{\rho)}h_M(\boldsymbol{\rho})(\omega_N)_{\alpha}.
\]
Then
\[
\mathcal{H}_{\boldsymbol{\rho}}\omega_M
= \sum_{N=0}^{\infty}\sum_{|\alpha|=N}(\mathcal{H}_{\boldsymbol{\rho}}\omega_M)_{\alpha}e_{\alpha}=
\sum_{N=0}^{\infty}
(-1)^{N+M}
(N+M)!h_N(\boldsymbol{\rho})h_M(\boldsymbol{\rho})
\omega_N.
\]
Equivalently, $U^{*}\mathcal{H}_{\boldsymbol{\rho}}U=\widetilde{B}_{\boldsymbol{\rho}}$, where
\[
\widetilde{B}_{\boldsymbol{\rho}}
=
\left(
(-1)^{N+M}(N+M)!h_N(\boldsymbol{\rho})h_M(\boldsymbol{\rho})
\right)_{N,M\geq 0}.
\]
Let $U_1:\ell^{2}(\mathbb{N}_{0})\to\ell^{2}(\mathbb{N}_{0})$
be the diagonal unitary $(U_1x)_N=(-1)^Nx_N.$ Then
\[
U_1\widetilde{B}_{\boldsymbol{\rho}}U_1
=
B_{\boldsymbol{\rho}},
\]
where
\[
(B_{\boldsymbol{\rho}})_{N,M}
=
(N+M)!h_N(\boldsymbol{\rho})h_M(\boldsymbol{\rho}).
\]
For $\boldsymbol{\rho}\neq \mathbf{0}$, observe that
\[
(N+M)!h_N(\boldsymbol{\rho})h_M(\boldsymbol{\rho})
=
c_N(\boldsymbol{\rho})c_M(\boldsymbol{\rho})
\binom{N+M}{N}\left(\frac{R_{\boldsymbol{\rho}}}{2}\right)^{N+M}.
\]
Hence
$B_{\boldsymbol{\rho}}
= D_{\boldsymbol{\rho}}H_{\frac{R_{\boldsymbol{\rho}}}{2}}D_{\boldsymbol{\rho}}.$

\end{proof}

\end{document}
